\section{Open questions (5)}

\subsection*{Q1. Non-Pauli noise: does the phase-diagram boundary survive?}
\textbf{Basis.} The paper's Fig.~6 and our 5-case reproduction both use 2-qubit
depolarizing noise on \texttt{CX} gates only. Real hardware exhibits amplitude
damping, dephasing, coherent over-rotation, cross-talk, and non-Markovian drift.
The mono-Exponential extrapolator's asymptote assumption (traceless observable
$\to 0$ under saturated depolarizing noise) is model-specific and may fail for
amplitude damping or coherent noise.

\textbf{Next steps.} Rerun the 5-case grid on Mitiq with (a) per-qubit
amplitude-damping + dephasing calibrated from an IBM device snapshot,
(b) a coherent-over-rotation channel, and (c) the paper's depolarizing model as
control. Compare winner-family boundaries. Add an asymptote-free
\texttt{ExpFactory} variant. Free stack (Aer + Open-Plan IBM Quantum).

\subsection*{Q2. Composition with virtual distillation and PEC}
\textbf{Basis.} The paper Sec.~V discusses composition with readout error
mitigation + Pauli twirling but does not treat virtual distillation or
probabilistic error cancellation (PEC). Whether pre-cleaned scans change the
extrapolation-family winner is unstudied.

\textbf{Next steps.} Prepend virtual distillation (Mitiq \texttt{ddd}, or a
hand-rolled $M$-copy purification) to the raw scan; rerun the 4-family ZNE
comparison. Repeat with PEC as pre-processor. Measure (i) whether Case D's
narrow-range Richardson blow-up is tempered by pre-distillation and (ii) whether
Case B's Exponential dominance survives.

\subsection*{Q3. Adaptive per-circuit best-practice selection}
\textbf{Basis.} The paper's Fig.~6 gives a static (depth, error-rate) lookup.
In practice, circuits with identical depth and $p_{2q}$ can have very different
observable-saturation behavior (brickwork vs QAOA vs QPE), so the family winner
may not be predictable from (depth, error) alone.

\textbf{Next steps.} Build a 20--30 circuit benchmark suite (brickwork, QAOA,
VQE ansatze, QPE, random Clifford+T). For each: run a 2-point pilot scan
($\lambda \in \{1, 1.5\}$, 500 shots), compute noisy-slope and
curvature-proxy, predict best family via a simple decision rule. Compare against
always-Linear, always-Exp, paper static recommendation, and oracle best.

\subsection*{Q4. Pulse-stretch vs folding: runtime overhead}
\textbf{Basis.} The paper (and our replication) uses statistical accuracy (RMSE)
as the only quality metric. Pulse stretching often requires OpenPulse access,
custom calibration, and can be slower per shot; digital folding is
$3$--$5\times$ more circuit depth for $\lambda \in \{3, 5\}$ but uses stock
transpilation.

\textbf{Next steps.} On IBM Quantum Open Plan (free), submit matched jobs:
digital fold at $\lambda \in \{1, 3, 5\}$ and pulse-stretched analog scaling at
the same effective $\lambda$. Measure wall-clock queue-to-result time,
per-shot execution time, and RMSE. Score by joint (RMSE, time-per-mitigated-obs).

\subsection*{Q5. ML-driven best-practice selector and cross-device transfer}
\textbf{Basis.} The paper offers a hand-crafted regime map. ML-driven QEM policy
selection is a natural extension but has not been evaluated for ZNE family
choice. Cross-device transfer (Falcon vs Eagle vs Heron) is the hard part.

\textbf{Next steps.} Generate 500--1000 (circuit, noise-model, best-family)
training tuples by running the 5-family bakeoff over random circuits and mixed
noise channels. Train an XGBoost classifier on 8 fingerprint features (2q depth,
T-count, entangling-gate density, ideal-obs magnitude, pilot noisy-obs,
pilot slope, pilot curvature, $p_{2q}$ estimate). Evaluate held-out RMSE vs
paper lookup and oracle. Free compute: uicgpu A100 for training, Aer for tuples.
